\chapter{Limitations}
\label{c:limitations}

The results reported in this thesis are obtained exclusively from synthetic
epidemic simulations governed by the SIR compartmental model. While the SIR
framework is a well-established abstraction for a broad class of infectious
diseases, it makes several simplifying assumptions that may not hold in
practice: infection probability is constant along each contact edge and does
not vary with dose or duration, every individual recovers with equal probability
per unit time, and there is no latency period between exposure and
infectiousness. Consequently, GNN models that learn to exploit the statistical
regularities of SIR dynamics---such as the characteristic shape of the
infected component at a given recovery ratio---may not generalise to outbreaks
driven by SEIR, SEIRS, or otherwise more complex dynamics. Performance on real
contact-tracing datasets, where the true source, spreading rate $\beta$, and
recovery rate $\mu$ are all unknown, could differ substantially from the
benchmarks reported here. Validation against real-world outbreaks with a known
index case, such as carefully documented laboratory or household transmission
studies, is therefore an important step that lies outside the scope of the
present work.

Scalability is a second significant limitation, affecting the two more
expressive temporal architectures in particular. The De Bruijn graph at order
$k$ has a node set whose size grows as $O(|E|^k)$ with the number of temporal
contacts $|E|$, and the De Bruijn GNN (DBGNN) \cite{Qarkaxhija2022} was
evaluated only on networks with at most a few thousand nodes and tens of
thousands of temporal contacts. The Temporal Event Graph (TEG) underlying the
DAG-GNN can contain $O(|E|^2)$ directed edges in the worst case, since every
pair of temporally adjacent contacts involving a shared node generates a link.
Neither architecture has been stress-tested on the large-scale contact networks
that arise in hospital surveillance or mobile-phone proximity studies, where
$|E|$ may reach millions. The static GNN and BN are considerably more
computationally tractable, but even these models face memory bottlenecks when
the number of Monte Carlo simulation runs used to construct node features is
large.

All four architectures share the single-source assumption: they assign a
probability to each node under the hypothesis that exactly one patient-zero
initiated the outbreak. This assumption simplifies training and evaluation but
is routinely violated in practice---many real-world epidemics are seeded by
multiple independent introduction events, particularly at the early stages of
an emerging pandemic. Models trained under the single-source assumption have no
mechanism to represent or detect such scenarios and will produce misleading
rankings when the true source set has cardinality greater than one. Finally,
while we follow the evaluation protocol recommended by \cite{Sterchi2025} to
mitigate inflated comparisons between learned methods and the Soft-Margin
Estimator baseline---which can achieve near-oracle performance when the
simulation model matches the data-generating process exactly---the comparison
remains constrained by the finite simulation budget, and differences between
architectures that are small in absolute terms should be interpreted with care.
